\documentclass{article}
\usepackage[utf8]{inputenc}
\usepackage[spanish]{babel}
\usepackage{listings}
\usepackage{graphicx}
\graphicspath{ {images/} }
\usepackage{cite}

\begin{document}

\begin{titlepage}
    \begin{center}
        \vspace*{1cm}
            
        \Huge
        \textbf{Proyecto de investigación}
            
        \vspace{1.5cm}
        \LARGE
        Taller de Memoria
            
        \vspace{1.5cm}
            
        \textbf{Erika Andrea Sanchez Salazar}
            
        \vfill
            
        \vspace{0.8cm}
            
        \Large
        Despartamento de Ingeniería Electrónica y Telecomunicaciones\\
        Universidad de Antioquia\\
        Medellín\\
        Septiembre de 2020
            
    \end{center}
\end{titlepage}

\tableofcontents
\newpage
\section{Introdución}\label{intro}
Esta es la primera sección, podemos agregar algunos elementos adicionales y todo será escrito correctamente. Más aún, si una palabra es demasiado larga y tiene que ser truncada, babel tratará de truncarla correctamente dependiendo del idioma.

\section{Taller de Memoria} \label{contenido}
Esta sección es para ver qué pasa con los comandos que definen texto.
\subsection{Defina que es la memoria del computador.}
la memoria principal del computador es un dispositivo de almacenamiento temporal y alta velocidad de acceso, que se utiliza para poder realizar las actividades mas rapido, ya que si se hiciera con discos duros seria mas lento todo el proceso.

Vamos a citar por ejemplo un artículo de \textbf{Albert Einstein} \cite{einstein}.
También es posible citar libros \cite{dirac} o documentos en línea \cite{knuthwebsite}.\\\\
Revisar en la última sección el formato de las referencias en IEEE.

\subsection{Mencione los tipos de memoria que conoce y haga una pequeña descripción de cada tipo.}
%
A continuación, se presenta el código 
\subsection{Describa la manera como se gestiona la memoria en un computador.}
En la sección \ref{imagenes}, se presentará como añadir ilustraciones al texto.
\subsection{¿Qué hace que una memoria sea más rápida que otra? ¿Por qué esto es importante?}

\section{Inclusión de imágenes} \label{imagenes}

En la Figura (\ref{fig:cpplogo}), se presenta el logo de C++ contenido en la carpeta images.

\begin{figure}[h]
\includegraphics[width=4cm]{cpplogo.png}
\centering
\caption{Logo de C++}
\label{fig:cpplogo}
\end{figure}

Las secciones (\ref{intro}), (\ref{contenido}) y (\ref{imagenes}) dependen del estilo del documento.

\bibliographystyle{IEEEtran}
\bibliography{references}

\end{document}
